% \iffalse meta-comment
% Copyright (C) 2026 Andreas Willibald Weber
%
% This work may be distributed and/or modified under the conditions of the
% LaTeX Project Public License, either version 1.3c of this license or (at your
% option) any later version. The latest version of the license is available at
% https://www.latex-project.org/lppl.txt.
%
% This work has the LPPL maintenance status `maintained'. The Current
% Maintainer of this work is Andreas Willibald Weber.
%
% This file is part of the `kanban' package; the complete list of files
% constituting the work appears in the main source file kanban.dtx.
% \fi
%
%    \begin{macrocode}
%<*package>
%    \end{macrocode}
%
% \subsection{Task collection}
%
% The task key family stores metadata without expanding the card text. A valid
% declaration appends one protected renderer call to the selected board-local
% list. Invalid or missing targets are diagnosed and never affect the counter.
% Collection and rendering are separate phases: this permits the environment
% to determine the tallest heading, the highest column, and the final column
% counts before any task card is typeset.
%
% \paragraph{Task-option keys.}
% The task family translates user-facing metadata into temporary task macros.
% Priority uses the validated semantic setter; author and deadline preserve
% their token lists; |due| aliases |deadline|; and |late| and |completed| are
% valueless flags. An unknown task option is ignored with a scoped warning.
%
% |\keys_define:nn| creates the |kanban/task| family. A |.code:n| handler sees
% the supplied value as |#1|. |.meta:n| redirects |due| through the canonical
% deadline handler. |.value_forbidden:n = true| tells |l3keys| that a Boolean
% flag must appear without |=...|, thereby distinguishing |completed| from a
% misleading form such as |completed=false|. The unknown handler uses
% |\l_keys_key_str| in the common named warning. The protected wrapper invokes
% |\keys_set:nn| and allows the surrounding LaTeX2e collector to remain outside
% LaTeX3 syntax.
%    \begin{macrocode}
\ExplSyntaxOn
\keys_define:nn {kanban/task}
  {
    prio .code:n = {\kanban@setpriority{#1}},
    author .code:n = {\def\kbtask@author{#1}},
    deadline .code:n = {\def\kbtask@deadline{#1}},
    due .meta:n = {deadline={#1}},
    late .code:n = {\def\kbtask@late{1}},
    late .value_forbidden:n = true,
    completed .code:n = {\def\kbtask@completed{1}},
    completed .value_forbidden:n = true,
    unknown .code:n =
      {\msg_warning:nnnn {kanban}{unknown-key}
        {task}{\l_keys_key_str}}
  }
\cs_new_protected:Npn \kanban@processtaskoptions #1
  {\keys_set:nn {kanban/task}{#1}}
\ExplSyntaxOff
%    \end{macrocode}
%
% \paragraph{Validation and deferred task collection.}
% The collector first gives every metadata field an empty value and processes
% the option list. It then requires an active board, a non-empty normalized
% target, and a target present in the board's identifier map. Each failed
% condition emits a specific error and ignores the task. A task inside a
% deliberately rejected nested board is silently discarded because the nesting
% error has already explained the cause; an ordinary task outside any board is
% diagnosed independently.
%
% |\ifbool| reads the |etoolbox| lifecycle flags. |\ifdefempty| tests the
% normalized column and identifier macros. |\PackageError| takes an error
% message and explanatory help; |\string| prints a command name without
% executing it, and |\MessageBreak| controls diagnostic layout. The nested
% branch structure ensures that neither counter nor list is modified until all
% target checks have succeeded.
%
% For a valid target, |\numgdef| globally defines the dynamic count macro to its
% current integer expression plus one. |\csxappto| globally appends expanded
% material to the dynamic task-list macro. Expansion is carefully selective:
% |\noexpand| retains the renderer command for the later rendering phase;
% |\csuse| resolves the already fixed column color; and |\unexpanded| protects
% the task body and all metadata from premature expansion. The
% |\expandafter| forms obtain the contents of temporary metadata macros before
% protecting those contents. The stored result is therefore a self-contained
% seven-argument call to |\kanban@rendercard| rather than a reference to task
% state that will be overwritten by the next declaration.
%    \begin{macrocode}
\newcommand\kanban@dotask[3][]{%
  \def\kbtask@prio{}\def\kbtask@author{}\def\kbtask@deadline{}%
  \def\kbtask@late{}\def\kbtask@completed{}%
  \kanban@processtaskoptions{#1}%
  \ifbool{kanban@active}
    {\kanban@normtaskcol{#2}%
     \ifdefempty\kbtask@column
       {\PackageError{kanban}
         {task has no assigned column; valid columns:
           \kanban@availablecolumns}
         {The first mandatory argument of
           \string\kanbantask\space must name a column from the
           enclosing board.\MessageBreak The task was ignored.}}
       {\ifdefempty\kbtask@columnid
       {\PackageError{kanban}
         {unknown kanban column '#2'; valid columns:
           \kanban@availablecolumns}
         {Column names are matched without regard to letter case.
           \MessageBreak
           The task was ignored because its normalized target is not in
           this board.}}
       {\expandafter\numgdef
           \csname kanban@count@\kanban@boardgroup
             @\kbtask@columnid\endcsname
           {\csname kanban@count@\kanban@boardgroup
             @\kbtask@columnid\endcsname+1}%
         \csxappto{kanban@list@\kanban@boardgroup
           @\kbtask@columnid}{%
           \noexpand\kanban@rendercard
           {\csuse{kanban@color@\kbtask@columnid}}{\unexpanded{#3}}%
           {\unexpanded\expandafter{\kbtask@prio}}%
           {\unexpanded\expandafter{\kbtask@author}}%
           {\unexpanded\expandafter{\kbtask@deadline}}%
           {\unexpanded\expandafter{\kbtask@late}}%
           {\unexpanded\expandafter{\kbtask@completed}}}%
       }}}
    {\ifbool{kanban@discardtasks}{}
      {\PackageError{kanban}
        {task declared outside a kanban environment}
        {Move the task declaration between \string\begin{kanban} and
          \string\end{kanban}. The task was ignored.}}}}
%    \end{macrocode}
% \begin{macro}{\kanbantask}
% \begin{macro}{\kbtask}
% The canonical public command supplies the compact task declaration syntax and
% delegates all key processing, validation, and collection to the internal
% implementation above. The concise command is installed by control-sequence
% equivalence only when its name is free; an existing definition is preserved
% with a warning.
%
% |\NewDocumentCommand| uses the signature |O{}mm|. |O{}| reads an optional
% bracketed key list and supplies an empty default when it is absent; the two
% |m| arguments read the target column and task body. The wrapper forwards the
% optional list in brackets because |\kanban@dotask| is a traditional LaTeX2e
% command with an optional argument. |\@ifundefined| performs the alias
% collision test, |\let| installs exact equivalence, and |\string| plus
% |\space| prints the affected command name legibly in the warning.
%    \begin{macrocode}
\NewDocumentCommand\kanbantask{O{}mm}{%
  \kanban@dotask[#1]{#2}{#3}}
\@ifundefined{kbtask}
  {\let\kbtask\kanbantask}
  {\PackageWarning{kanban}
    {\string\kbtask\space is already defined;
      \MessageBreak use \string\kanbantask\space instead}}
%    \end{macrocode}
% \end{macro}
% \end{macro}
%
% \subsection{Board lifecycle}
%
% Board keys update only the active board snapshot. Column setup allocates its
% isolated lists, prepares both rendering sequences, and validates the complete
% gap width before typesetting begins. The entry routine also rejects nesting
% before any outer state can be replaced. The environment itself supplies a TeX
% group, so assignments intentionally divided into persistent and local forms
% acquire the documented lifetime without a separate restoration procedure.
%
% \paragraph{Board column selection and board-option keys.}
% A blank board column value means ``use the persistent default'', while a
% non-blank value is normalized as a board-local override. The board key family
% then provides local counterparts to the persistent setup keys and adds the
% valueless |break| switch. Start and reference dates use the same deferred
% strategy as persistent setup so key order cannot alter ambiguous-date parsing.
%
% |\ifblank| selects default initialization or validated local selection.
% |\keys_define:nn| creates the |kanban/board| family. The |break| handler calls
% |\booltrue|, and |.value_forbidden:n| requires flag syntax without a value.
% |.code:n| handlers delegate to the board-local selectors; |.meta:n| provides
% the documented spelling aliases; and the date handlers stage values with
% |\tl_set:Nn| and |\bool_set_true:N|. The nested |labels| key invokes the
% board-label family. As in the other key families, the |unknown| handler uses
% |\l_keys_key_str| to identify the ignored option precisely.
%    \begin{macrocode}
\newcommand\kanban@setboardcolumns[1]{%
  \ifblank{#1}
    {\kanban@columnsboardinit}
    {\kanban@selectboardcolumns{#1}}}
\ExplSyntaxOn
\keys_define:nn {kanban/board}
  {
    break .code:n = {\booltrue{kanban@break}},
    break .value_forbidden:n = true,
    columns .code:n = {\kanban@setboardcolumns{#1}},
    date-format .code:n = {\kanban@selectboarddateformat{#1}},
    dateformat .meta:n = {date-format={#1}},
    startdate .code:n =
      {
        \tl_set:Nn \l__kanban_pending_startdate_tl {#1}
        \bool_set_true:N \l__kanban_pending_startdate_bool
      },
    reference-date .code:n =
      {
        \tl_set:Nn \l__kanban_pending_reference_date_tl {#1}
        \bool_set_true:N \l__kanban_pending_reference_date_bool
      },
    today .meta:n = {reference-date={#1}},
    palette .code:n = {\kanban@selectboardpalette{#1}},
    mode .code:n = {\kanban@selectboardmode{#1}},
    columnsep .code:n = {\kanban@setboardcolumnsep{#1}},
    labels .code:n = {\keys_set:nn {kanban/board-labels}{#1}},
    unknown .code:n =
      {\msg_warning:nnnn {kanban}{unknown-key}
        {board}{\l_keys_key_str}}
  }
%    \end{macrocode}
%
% \paragraph{A deferred board-option transaction.}
% Board-option processing mirrors persistent setup but commits staged dates to
% local snapshot setters. All immediate options, including a possible
% |date-format| override, are therefore resolved before the dates are parsed.
% The transaction runs after the persistent values have been copied into the
% board, so omitted keys naturally inherit their defaults.
%
% |\bool_set_false:N| resets the two staging sentinels, |\keys_set:nn| processes
% the complete option list, and |\bool_if:NT| performs only the requested
% commits. |\exp_args:NV| supplies each stored token-list value to its setter.
% The protected definition prevents assignment code from being expanded in a
% moving context. |\ExplSyntaxOff| returns to the LaTeX2e syntax used by the
% dynamic renderer-sequence construction below.
%    \begin{macrocode}
\cs_new_protected:Npn \kanban@processboardoptions #1
  {
    \bool_set_false:N \l__kanban_pending_startdate_bool
    \bool_set_false:N \l__kanban_pending_reference_date_bool
    \keys_set:nn {kanban/board}{#1}
    \bool_if:NT \l__kanban_pending_startdate_bool
      {\exp_args:NV \kanban@setboardstartdate
        \l__kanban_pending_startdate_tl}
    \bool_if:NT \l__kanban_pending_reference_date_bool
      {\exp_args:NV \kanban@setboardreferencedate
        \l__kanban_pending_reference_date_tl}
  }
\ExplSyntaxOff
%    \end{macrocode}
%
% \paragraph{Per-column state and renderer-sequence construction.}
% Each normalized column is represented internally by a numeric identifier.
% Setup stores its display name, initializes its board-qualified deferred task
% list and counter, assigns the cyclic palette background and foreground, and
% appends one adapter call to each renderer sequence. Board qualification is
% used for card lists and counts because those definitions are global during
% collection; simple color and name definitions are removed explicitly when
% the board closes.
%
% |\csdef| and |\csgdef| from |etoolbox| define dynamically named local and
% global control sequences. |\csedef| performs a local expanded definition, so
% the selected colors are fixed immediately. |\numdef| evaluates and stores the
% incremented number of columns. |\eappto| appends expanded material to the
% header, fixed-body, and break-layout sequences, while |\noexpand| preserves
% the adapter command for execution at rendering time. |\ifnumequal| treats the
% first break column specially; later columns prepend |\switchcolumn| so
% |paracol| receives one stream per board column.
%    \begin{macrocode}
\newcommand\kanban@setupcolumn[2]{%
  \csdef{kanban@name@#2}{#1}%
  \csgdef{kanban@list@\kanban@boardgroup @#2}{}%
  \csgdef{kanban@count@\kanban@boardgroup @#2}{0}%
  \kanban@palettecolor{#2}\csedef{kanban@color@#2}{\kb@thiscolor}%
  \csedef{kanban@headingtext@#2}{\kb@thisheadingtext}%
  \numdef\kanban@ncols{\kanban@ncols+1}%
  \eappto\kanban@headseq{\noexpand\kanban@onehead{#2}}%
  \eappto\kanban@colseq{\noexpand\kanban@onecol{#2}}%
  \ifnumequal{#2}{1}
    {\eappto\kanban@colseq@b{\noexpand\kanban@onecol@b{#2}}}
    {\eappto\kanban@colseq@b{\noexpand\switchcolumn
       \noexpand\kanban@onecol@b{#2}}}}
%    \end{macrocode}
%
% \paragraph{Whole-board geometry validation.}
% A syntactically valid non-negative gap may still be impossible for a board
% with many columns. Before fixed rendering, the validator computes the total
% width of the $N-1$ inter-column gaps and requires that total to be strictly
% smaller than |\linewidth|. An impossible value produces an actionable error
% and falls back to zero for this board, allowing the document to continue.
%
% The primitive |\ifnum| skips multiplication for a one-column board.
% Dimension-register assignment copies the accepted board gap to
% |\kanban@totalgaps|, and |\multiply| scales it by the integer expression
% |\numexpr\kanban@ncols-1\relax|. |\ifdim| compares the result with the
% available line width. Its empty true branch accepts valid geometry; the false
% branch calls |\PackageError| and locally assigns |0pt|. Each primitive
% conditional is terminated by its matching |\fi|.
%    \begin{macrocode}
\newcommand\kanban@validateboardgeometry{%
  \ifnum\kanban@ncols>1\relax
    \kanban@totalgaps=\kanban@boardcolumnsep
    \multiply\kanban@totalgaps by \numexpr\kanban@ncols-1\relax
    \ifdim\kanban@totalgaps<\linewidth\relax\else
      \PackageError{kanban}
        {column separation leaves no width for columns}
        {For N columns, (N-1) times columnsep must be smaller than
          \string\linewidth.\MessageBreak
          The current board will use columnsep=0pt.}%
      \kanban@boardcolumnsep=0pt
    \fi
  \fi}
%    \end{macrocode}
%
% \paragraph{Environment entry and state snapshot.}
% Opening a board must either reject unsupported nesting or create a complete,
% isolated snapshot before tasks are read. The successful path resets lifecycle
% flags, advances the unique board identifier, copies every persistent scalar
% and list default, processes local overrides, materializes the visual mode,
% clears renderer sequences, registers all columns, validates aggregate
% geometry, and suppresses source whitespace before the first task.
%
% The initial |\ifbool| detects an already active board before any outer state
% is overwritten. A rejected inner environment clears its render and active
% flags locally, sets |kanban@discardtasks|, and issues one |\PackageError|; its
% body can then be read safely without being attached to the outer board. On a
% valid entry, |\stepcounter| advances the LaTeX counter and |\arabic| supplies
% its printable value. |\edef| takes expanded scalar snapshots, while |\let|
% copies the current label meanings. Length-register assignment copies the
% persistent gap. The previously defined initialization and option functions
% then build the structured state. Finally, |\def| empties the three deferred
% sequences and resets the numeric column macro, and |\ignorespaces| removes
% indentation introduced by the source layout immediately after
% |\begin{kanban}|.
%    \begin{macrocode}
\newcommand\kanban@begin[1]{%
   \ifbool{kanban@active}
     {\boolfalse{kanban@render}%
      \boolfalse{kanban@active}%
      \booltrue{kanban@discardtasks}%
      \PackageError{kanban}
        {nested kanban environments are not supported}
        {Close the current board before opening another one.
          \MessageBreak The nested board was ignored.}}
     {\booltrue{kanban@render}%
   \boolfalse{kanban@break}%
   \boolfalse{kanban@discardtasks}%
   \booltrue{kanban@active}%
   \stepcounter{kanban@board}%
   \edef\kanban@boardgroup{\arabic{kanban@board}}%
   \edef\kanban@thestart{\kanban@startdate}%
   \edef\kanban@boarddateformat{\kanban@dateformat}%
   \edef\kanban@thereference{\kanban@referencedate}%
   \edef\kanban@boardmode{\kanban@mode}%
   \let\kanban@boardhighlabel\kanban@highlabel
   \let\kanban@boardmediumlabel\kanban@mediumlabel
   \let\kanban@boardlowlabel\kanban@lowlabel
   \let\kanban@boarddonelabel\kanban@donelabel
   \kanban@boardcolumnsep=\kanban@columnsep
   \kanban@paletteboardinit
   \kanban@columnsboardinit
   \kanban@processboardoptions{#1}%
   \kanban@applymode
   \def\kanban@ncols{0}\def\kanban@headseq{}%
   \def\kanban@colseq{}\def\kanban@colseq@b{}%
   \kanban@setupboardcolumns
   \kanban@validateboardgeometry
   \ignorespaces}}
%    \end{macrocode}
% \begin{macro}{\kanban}
% \begin{macro}{\endkanban}
% The public environment has one optional board-key list and no mandatory
% column argument; columns come from the current persistent default unless a
% board-local |columns| key overrides it. Closing the environment selects the
% page-breakable |paracol| renderer or the indivisible fixed raster. A rejected
% nested environment has its render flag cleared and therefore produces no
% board. Successfully rendered boards release their dynamically named state
% afterward.
%
% |\NewDocumentEnvironment| defines both the begin and end commands and checks
% that the environment name is new. Its signature |O{}| reads an optional key
% list with an empty default. The begin code delegates to |\kanban@begin|. In
% the end code, |\ifbool{kanban@render}| guards the complete rendering and
% cleanup path. |\par\medskip| leaves the task-declaration paragraph and
% separates the board from preceding text.
%
% In break mode, |\setlength| assigns the board gap to |paracol|'s standard
% |\columnsep|, |\begin{paracol}{\kanban@ncols}| opens the requested number of
% independent streams, and the prepared sequence supplies headers, cards, and
% |\switchcolumn| calls. In fixed mode, |\noindent| suppresses paragraph
% indentation and a full-line |minipage| makes the board indivisible. The
% |tcbraster| has one column per Kanban stage and two rows: all headers followed
% by all bodies. Its keys set the available width, board gap, zero row gap,
% equal height within each row, and top alignment. Executing
% |\kanban@headseq\kanban@colseq| fills the raster in row-major order.
% After either renderer, |\vspace{2mm}| provides documented following space and
% |\kanban@cleanupboard| undefines per-column state. The environment's ordinary
% TeX group then restores all remaining local snapshots automatically.
%    \begin{macrocode}
\NewDocumentEnvironment{kanban}{O{}}
  {\kanban@begin{#1}}
  {\ifbool{kanban@render}{\par\medskip
   \ifbool{kanban@break}
     {\setlength\columnsep{\kanban@boardcolumnsep}%
      \begin{paracol}{\kanban@ncols}\kanban@colseq@b\end{paracol}}
     {\noindent\begin{minipage}{\linewidth}%
      \begin{tcbraster}[
        raster columns=\kanban@ncols,
        raster rows=2,
        raster width=\linewidth,
        raster column skip=\kanban@boardcolumnsep,
        raster row skip=0pt,
        raster before skip=0pt,
        raster after skip=0pt,
        raster equal height=rows,
        raster valign=top]
        \kanban@headseq\kanban@colseq
      \end{tcbraster}%
      \end{minipage}}%
   \par\vspace{2mm}\kanban@cleanupboard}{}}
%</package>
%    \end{macrocode}
% \end{macro}
% \end{macro}
% \Finale
